% VI23_FULL_SOLUTION_REFINEMENT_V1
\subsection*{VI.23.P10: Pullback of a principal open}
\begin{problem}\label{prob:vi23-p10}
Let $A\to B$ and $f\in A$. Compute
\[
\Spec B\times_{\Spec A}\Spec A_f.
\]
\end{problem}

\ifdefined\FullProblemDossiers
\paragraph{Why this problem matters.}
This shows concretely that inverse images of principal opens are principal opens.

\paragraph{Useful ideas before starting.}
Translate the requested statement into compatible maps from an arbitrary test scheme, and use uniqueness before choosing coordinates.

\paragraph{Guided hints.}
\begin{enumerate}
\item Use $B\otimes_AA_f$.
\item Show the image of $f$ becomes invertible.
\item Invoke the localization universal property twice.
\end{enumerate}

\begin{solution}
Let $\varphi:A\to B$.  By the affine theorem the coordinate ring is
\[
B\otimes_AA_f.
\]
The map $B\to B_{\varphi(f)}$ together with the localization map
$A_f\to B_{\varphi(f)}$ gives an $A$-balanced multiplication map
\[
B\otimes_AA_f\to B_{\varphi(f)},\qquad
b\otimes\frac{a}{f^n}\mapsto\frac{b\varphi(a)}{\varphi(f)^n}.
\]
Conversely the canonical map $B\to B\otimes_AA_f$ sends $\varphi(f)$ to the unit
$1\otimes f$, so it factors uniquely through $B_{\varphi(f)}$.  The two maps are inverse by the
universal property of localization.

Hence
\[
\Spec B\times_{\Spec A}\Spec A_f
\cong\Spec B_{\varphi(f)}
=D(\varphi(f))\subseteq\Spec B.
\]
The projection to $\Spec B$ is precisely the standard open immersion.
\end{solution}

\paragraph{What happened mathematically.}
A geometric assertion was reduced to a representability statement, so the canonical map was forced by universality.

\paragraph{Extensions.}
Repeat the argument after restricting to affine opens and compare with the tensor-product formula.

\paragraph{Provenance.}
\textbf{Canonical fiber-product computation; systematic base-change language is reserved for VI/24.}
\else
\paragraph{Provenance.} \textbf{Canonical fiber-product computation; systematic base-change language is reserved for VI/24.}
\fi
